\DocumentMetadata{
  pdfversion=2.0,
  pdfstandard=ua-2,
  lang=en-US,
  tagging=on,
  tagging-setup={math/setup=mathml-SE,tabsorder=structure}
}
\documentclass[11pt,letterpaper]{article}
\usepackage[margin=0.68in,includefoot]{geometry}
\usepackage{newcomputermodern}
\usepackage{amsmath,amscd,mathtools}
\usepackage{tikz,adjustbox}
\providecommand{\square}{\mdlgwhtsquare}
\usepackage[hidelinks]{hyperref}
\hypersetup{
  pdftitle={Numerical Analysis PhD Exam, January 2011},
  pdfauthor={University of Florida Department of Mathematics},
  pdfsubject={Graduate mathematics examination},
  pdfdisplaydoctitle=true,
  bookmarksopen=true
}
\setcounter{secnumdepth}{0}
\setlength{\parindent}{0pt}
\setlength{\parskip}{0.28em}
\setlength{\emergencystretch}{3em}
\setlength{\leftmargini}{2.15em}
\setlength{\leftmarginii}{2.65em}
\setlength{\leftmarginiii}{3em}
\renewcommand{\labelenumi}{\textbf{\theenumi.}}
\pagestyle{empty}
\raggedbottom
\allowdisplaybreaks
\newenvironment{examproblems}[1]{%
  \begin{enumerate}%
  \setcounter{enumi}{#1}%
  \setlength{\topsep}{0.35em}%
  \setlength{\partopsep}{0pt}%
  \setlength{\itemsep}{0.48em}%
  \setlength{\parsep}{0.18em}%
}{\end{enumerate}}
\newenvironment{examsubparts}{%
  \begin{enumerate}%
  \setlength{\topsep}{0.18em}%
  \setlength{\partopsep}{0pt}%
  \setlength{\itemsep}{0.16em}%
  \setlength{\parsep}{0.1em}%
}{\end{enumerate}}
\newenvironment{examinstructions}{%
  \par\begingroup\itshape
}{%
  \par\endgroup\vspace{0.18em}
}
\makeatletter
\renewcommand\section{\@startsection{section}{1}{0pt}%
  {-1.25ex plus -.25ex minus -.1ex}%
  {0.55ex plus .1ex}%
  {\normalfont\large\bfseries}}
\renewcommand{\@maketitle}{%
  \newpage\null\vskip -1em%
  {\footnotesize\bfseries
    \href{https://www.ufl.edu/}{University of Florida}, \href{https://math.ufl.edu/}{Department of Mathematics}\par}%
  \vskip -0.5em%
  \tagstructbegin{tag=Title}%
    {\LARGE\bfseries\href{https://gma.math.ufl.edu/exams/numerical-analysis-phd/}{Numerical Analysis PhD Exam}, January 2011\par}%
  \tagstructend%
  \vskip 0.25em%
  \tagmcbegin{artifact}%
    \hrule height 1pt%
  \tagmcend%
  \vskip 1em%
}
\makeatother
\title{Numerical Analysis PhD Exam, January 2011}
\author{}
\date{}
\begin{document}
\maketitle
\thispagestyle{empty}
\begin{examinstructions}
Answer any 8 questions.
\end{examinstructions}
\begin{examproblems}{0}
\item
Suppose~\(A\) and~\(B\) are square matrices such that \(AB\) is normal. Prove that \(\|AB\|_2 \leq \|BA\|\). (We use \(\|\cdot\|_2\) to denote the spectral norm and \(\|\cdot\|\) to denote any induced matrix norm.)
\item
If~\(P\) is a projector that is neither~\(0\) nor the identity, prove that \(\|P\|=\|I-P\|\) in any norm induced by a vector norm generated by an inner product.
\item
Suppose~\(A\) is a Hermitian positive definite matrix split into \(A=C+C^*+D\) where~\(D\) is also Hermitian positive definite. Prove that \(B=C+\omega^{-1}D\) is invertible whenever \(0<\omega<2\). Consider the iteration \(x_{n+1}=x_n+B^{-1}(b-Ax_n)\), with any initial iterate \(x_0\). Prove that \(x_n\) converges to \(x=A^{-1}b\) whenever \(0<\omega<2\).
\item
If a square matrix~\(A\) can be block partitioned as \(A=\begin{bmatrix}0&M\\N&0\end{bmatrix}\), prove that \(\sigma(A)=\{z\in\mathbb{C}:z^2\in\sigma(MN)\cup\sigma(NM)\}\).
\item
Let \(A\in\mathbb{R}^{N\times N}\) and \(b\in\mathbb{R}^N\). Consider the following iteration for solving \(Ax=b\), that computes \(x_{n+1}\), given \(x_n\in\mathbb{R}^N\), as follows (in~\(m\) intermediate steps): Setting \(x_{n+1}^{(0)}=x_n\), for \(\ell=1,2,\ldots,m\), compute \(x_{n+1}^{(\ell)}=x_{n+1}^{(\ell-1)}+\tau_\ell(b-Ax_{n+1}^{(\ell-1)})\). Then, define \(x_{n+1}=x_{n+1}^{(m)}\). There is a linear operator~\(E\) such that \(x_{n+1}-x=E(x_n-x)\). Give a formula for~\(E\). Suppose~\(A\) is Hermitian and positive definite with spectral condition number~\(\kappa\). Prove that there are real values of the~\(m\) parameters \(\tau_\ell\) such that 
\[
\rho(E)\leq 2\left(\frac{\sqrt{\kappa}-1}{\sqrt{\kappa}+1}\right)^m.
\]
\item
(The Aitken accelerated convergence algorithm). Let \(\alpha=\phi(\alpha)\) be a fixed point of \(\phi(x)\). Suppose that~\(\phi\) is twice continuously differentiable with \(\phi'(\alpha)\neq1\). Consider the iteration 
\[
x_{n+1}=\Phi(x_n)
\]
 
\[
\Phi(x)=\phi(\phi(x))+\frac{H(x)}{1-H(x)}\left(\phi(\phi(x))-\phi(x)\right)
\]
 
\[
H(x)=\frac{\phi(\phi(x))-\phi(x)}{\phi(x)-x}.
\]
 Show that the iteration \(x_{n+1}=\Phi(x_n)\) is quadratically convergent to~\(\alpha\) when the starting guess \(x_0\) is sufficiently close to~\(\alpha\).
\item
Let~\(n\) be a positive integer, \(h=1/n\), and consider the grid of points \((ih,jh)\) for \(i,j=0,1,\ldots,n\). Let~\(A\) be the finite difference operator with the “5-point stencil” discretizing the Laplace operator \(-\Delta\), with zero Dirichlet boundary conditions on this grid. Describe it. Prove that the spectrum of~\(A\) consists of the numbers 
\[
\lambda_{lm}=4h^{-2}\left(\sin^2(l\pi h/2)+\sin^2(m\pi h/2)\right),
\]
 for all \(l,m=1,\ldots,n-1\).
\item
Let \(S_n^1\) denote the space of linear splines based on knots \(a=t_0<t_1<\cdots<t_n=b\), i.e., \(S_n^1=\{v:v|_{[t_j,t_{j+1}]}\text{ is linear for all }j=0,1,\ldots,n-1\text{ and }v\text{ is continuous on }[a,b]\}\). Let \(s_f\in S_n^1\) interpolate a continuous function~\(f\) at the knots. Prove that \(\|f-s_f\|_\infty\leq2\|f-s\|_\infty\) for any \(s\in S_n^1\).
\item
Given that the Newton iteration for finding a root~\(r\) of~\(f\) converges, \(f\in C^3(\mathbb{R})\), and \(f(r)=f'(r)=0\neq f''(r)\), prove that the convergence cannot be quadratic. Suggest a modification that restores quadratic convergence for smooth~\(f\).
\item
Let \(x_0,x_1,\ldots,x_n\) be distinct numbers and for each~\(i\), let \(\ell_i\) denote the product of all \((x_i-x_j)^{-1}\) for all \(j\neq i\), \(j=0,\ldots,n\). Express the number \(q_f=f(x_0)\ell_0+f(x_1)\ell_1+\cdots+f(x_n)\ell_n\) as a divided difference of~\(f\). Then show that \(q_f=0\) whenever~\(f\) is a polynomial of degree \(n-1\).
\end{examproblems}
\end{document}
